\section{讨论与有效性威胁}
\subsection{哪些结论目前可信}
版本与来源窗口核验、受限参数契约和明确构造任务的验收，具有可直接复查的内部依据。
其中的可信性来自目标、输入、参照和反例的明确性，而不是模型品牌或论文引用数量。
三种实现路线对24个CSV参考结果一致，390项验收器诊断及代表案例手算核对，
降低了把一处实现错误误认成模型失败的风险；但它们仍共享研究者规定的任务语义。
本文没有独立人类质量标签，不能估计八维评分的MAE、人工一致性或语义支持准确率。

新产物实验的优势是把真实模型输出与独立目标连接起来，提供可反驳的失败，而不是由模型自报成功。
其限制也很明确：这是一个文本执行器，不使用Python或社区脚本。
原始技能建议运行工具，而当前条件不允许；因此结果刻画的是该工具受限条件下的文档注入效用，
不证明原技能在正常工具环境中同样低效。
一个确定性CSV求解器可以直接完成这些刚性任务，后续应将其作为工具路线的强对照，
而不能把“用确定性工具解决计数”包装成新颖算法。

\subsection{组合方法的增量与尚未成立的优势}
本协议把前人的多维评审、原子证据、可执行验收与集合选择连接为一个可审计流程。
实证支持的增量是：发现某些卡片参数名检查会漏掉的机械错误，
把模型提出的引用绑定到受控来源，并避免把范围外处置与完整产物成功混合报告。
标准库强对照同样检出构造错误，原包加契约未建立稳定产物优势，网络代理重排也没有普遍改善。
这些负结果限定了当前主张，不应被“更多维度”“更多反思”或“更复杂网络”取代。

仅契约组在本开发设计中使用较少上下文，并取得最高的观察成功数。
这提示后续可研究面向任务的精炼操作约束，但它不是证明简短技能总是更好。
需要扩大目标技能数、任务类型和模型，并在结果产生前固定比较与预算。
若最终只发现便宜的确定性工具已足够，应优先采用工具，而不是为了展示智能体复杂性增加开销。

\subsection{来源、任务与模型的外推限制}
D1只覆盖两个来源；原始1,881对象也不是1,881个独立仓库。
补充冻结的51个外部候选文档有5个与旧内容连接，按原仓库来源与近重关系联合分组规则，
两个新仓库都与旧来源相连，不能直接得到独立外部测试。
现有D2中的卡片混入问题仍需修正；数量凑足也不等于代表性足够。

12个新构造条件刻意包含CSV和日期陷阱，不代表真实用户任务的发生频率。
它们是开发条件，答案由研究者构造并对模型隐藏，不是独立专家认可的公开基准。
五臂、两次重复不能当作300个独立任务。
单一模型别名可能随服务更新；关闭思考和没有工具也限制结论。
因此，本文的条件区间描述该设计中的变动，不构成跨来源、跨模型或未来请求的保证。

\subsection{下一阶段的可检验路线}
第一，建立真实双盲语义标注和仲裁，比较发表版SkillNet、SkillCT画像估计、
同材料普通Judge及证据约束方法，并披露仲裁前一致性与未知比例。
第二，在合格隔离环境中允许真实工具，使用新的输入和独立产物验收，
比较无技能、原始包、精炼契约和确定性工具路线。
第三，只有单技能证据的效度得到验证后，才把它作为固定输入研究条件收益、
互补集合和验证预算分配，并检查网络更新后的旧任务遗忘\citep{P03,P06,P07,P23}。
当前研究尚不能给出经过校准的收益下界，故没有给未知节点填入任意“可信概率”。

样本量也应服从可检出效应，而不是整齐的任务数量。
假定配对不一致概率20\%、双侧阈值0.05，100个独立二元单位检测5个百分点提升的功效约14.3\%，
600个约76.0\%。
这只是规划敏感性，不是本研究观察到的效应或复杂任务族的直接功效保证。
后续需要用开发数据估计不一致率与相关性，再冻结正式留出规模。
